\section{Closed subschemes of relative proj}
\label{section-closed-in-relative-proj}

\noindent
Some auxiliary lemmas about closed subschemes of relative proj.

\begin{lemma}
\label{lemma-closed-subscheme-proj}
Let $S$ be a scheme. Let $\mathcal{A}$ be a quasi-coherent graded
$\mathcal{O}_S$-algebra. Let
$p : X = \underline{\text{Proj}}_S(\mathcal{A}) \to S$ be the relative
Proj of $\mathcal{A}$. Let $i : Z \to X$ be a closed subscheme. Denote
$\mathcal{I} \subset \mathcal{A}$ the kernel of the canonical map
$$
\mathcal{A}
\longrightarrow
\bigoplus\nolimits_{d \geq 0} p_*\left((i_*\mathcal{O}_Z)(d)\right).
$$
If $p$ is quasi-compact, then there is an isomorphism
$Z = \underline{\text{Proj}}_S(\mathcal{A}/\mathcal{I})$.
\end{lemma}

\begin{proof}
The morphism $p$ is separated by
Constructions, Lemma \ref{constructions-lemma-relative-proj-separated}.
As $p$ is quasi-compact, $p_*$ transforms quasi-coherent modules
into quasi-coherent modules, see
Schemes, Lemma \ref{schemes-lemma-push-forward-quasi-coherent}.
Hence $\mathcal{I}$ is a quasi-coherent $\mathcal{O}_S$-module.
In particular, $\mathcal{B} = \mathcal{A}/\mathcal{I}$ is a
quasi-coherent graded $\mathcal{O}_S$-algebra. The functoriality
morphism $Z' = \underline{\text{Proj}}_S(\mathcal{B}) \to
\underline{\text{Proj}}_S(\mathcal{A})$ is everywhere defined and
a closed immersion, see Constructions, Lemma
\ref{constructions-lemma-surjective-graded-rings-map-relative-proj}.
Hence it suffices to prove $Z = Z'$ as closed subschemes of $X$.

\medskip\noindent
Having said this, the question is local on the base and we may assume
that $S = \Spec(R)$ and that $X = \text{Proj}(A)$ for some
graded $R$-algebra $A$. Assume $\mathcal{I} = \widetilde{I}$
for $I \subset A$ a graded ideal. By
Constructions, Lemma \ref{constructions-lemma-proj-quasi-compact}
there exist $f_0, \ldots, f_n \in A_{+}$ such that
$A_{+} \subset \sqrt{(f_0, \ldots, f_n)}$ in other words
$X = \bigcup D_{+}(f_i)$. Therefore, it suffices to check that
$Z \cap D_{+}(f_i) = Z' \cap D_{+}(f_i)$ for each $i$.
By renumbering we may assume $i = 0$.
Say $Z \cap D_{+}(f_0)$, resp.\ $Z' \cap D_{+}(f_0)$
is cut out by the ideal $J$, resp.\ $J'$ of $A_{(f_0)}$.

\medskip\noindent
The inclusion $J' \subset J$.
Let $d$ be the least common multiple of $\deg(f_0), \ldots, \deg(f_n)$.
Note that each of the twists $\mathcal{O}_X(nd)$ is invertible, trivialized
by $f_i^{nd/\deg(f_i)}$ over $D_{+}(f_i)$, and that for any quasi-coherent
module $\mathcal{F}$ on $X$ the multiplication maps
$\mathcal{O}_X(nd) \otimes_{\mathcal{O}_X} \mathcal{F}(m)
\to \mathcal{F}(nd + m)$ are isomorphisms, see
Constructions, Lemma \ref{constructions-lemma-when-invertible}.
Observe that $J'$ is the ideal generated by the elements $g/f_0^e$ where
$g \in I$ is homogeneous of degree $e\deg(f_0)$ (see proof of
Constructions, Lemma
\ref{constructions-lemma-surjective-graded-rings-map-proj}).
Of course, by replacing $g$ by $f_0^lg$ for suitable $l$
we may always assume that $d | e$. Then, since $g$ vanishes as a section of
$\mathcal{O}_X(e\deg(f_0))$ restricted to $Z$ we see that
$g/f_0^d$ is an element of $J$. Thus $J' \subset J$.

\medskip\noindent
Conversely, suppose that $g/f_0^e \in J$. Again we may assume $d | e$.
Pick $i \in \{1, \ldots, n\}$. Then $Z \cap D_{+}(f_i)$ is
cut out by some ideal $J_i \subset A_{(f_i)}$. Moreover,
$$
J \cdot A_{(f_0f_i)} = J_i \cdot A_{(f_0f_i)}.
$$
The right hand side is the localization of $J_i$ with respect to
$f_0^{\deg(f_i)}/f_i^{\deg(f_0)}$. It follows that
$$
f_0^{e_i}g/f_i^{(e_i + e)\deg(f_0)/\deg(f_i)} \in J_i
$$
for some $e_i \gg 0$ sufficiently divisible. This proves that
$f_0^{\max(e_i)}g$ is an element of $I$, because its restriction to each
affine open $D_{+}(f_i)$ vanishes on the closed subscheme
$Z \cap D_{+}(f_i)$. Hence $g/f_0^e \in J'$ and we conclude $J \subset J'$
as desired.
\end{proof}

\begin{example}
\label{example-closed-subscheme-of-proj}
Let $A$ be a graded ring. Let $X = \text{Proj}(A)$ and $S = \Spec(A_0)$.
Given a graded ideal $I \subset A$ we obtain a closed subscheme
$V_+(I) = \text{Proj}(A/I) \to X$ by Constructions, Lemma
\ref{constructions-lemma-surjective-graded-rings-map-proj}.
Translating the result of Lemma \ref{lemma-closed-subscheme-proj}
we see that if $X$ is quasi-compact, then any closed subscheme $Z$
is of the form $V_+(I(Z))$ where the graded ideal $I(Z) \subset A$
is given by the rule
$$
I(Z) = \Ker(A \longrightarrow
\bigoplus\nolimits_{n \geq 0} \Gamma(Z, \mathcal{O}_Z(n)))
$$
Then we can ask the following two natural questions:
\begin{enumerate}
\item Which ideals $I$ are of the form $I(Z)$?
\item Can we describe the operation $I \mapsto I(V_+(I))$?
\end{enumerate}
We will answer this when $A$ is Noetherian.

\medskip\noindent
First, assume that $A$ is generated by $A_1$ over $A_0$. In this case,
for any ideal $I \subset A$ the kernel of the map
$A/I \to \bigoplus \Gamma(\text{Proj}(A/I), \mathcal{O})$
is the set of torsion elements of $A/I$, see
Cohomology of Schemes, Proposition
\ref{coherent-proposition-coherent-modules-on-proj}.
Hence we conclude that
$$
I(V_+(I)) = \{x \in A \mid A_n x \subset I\text{ for some }n \geq 0\}
$$
The ideal on the right is sometimes called the saturation of $I$.
This answers (2) and the answer to (1) is that an ideal is
of the form $I(Z)$ if and only if it is saturated, i.e., equal
to its own saturation.

\medskip\noindent
If $A$ is a general Noetherian graded ring, then we use
Cohomology of Schemes, Proposition
\ref{coherent-proposition-coherent-modules-on-proj-general}.
Thus we see that for $d$ equal to the lcm of the degrees
of generators of $A$ over $A_0$ we get
$$
I(V_+(I)) = \{x \in A \mid (Ax)_{nd} \subset I\text{ for all }n \gg 0\}
$$
This can be different from the saturation of $I$ if $d \not = 1$.
For example, suppose that $A = \mathbf{Q}[x, y]$
with $\deg(x) = 2$ and $\deg(y) = 3$. Then $d = 6$.
Let $I = (y^2)$. Then we see $y \in I(V_+(I))$ because
for any homogeneous $f \in A$ such that $6 | \deg(fy)$
we have $y | f$, hence $fy \in I$. It follows that
$I(V_+(I)) = (y)$ but $x^n y \not \in I$ for all $n$
hence $I(V_+(I))$ is not equal to the saturation.
\end{example}

\begin{lemma}
\label{lemma-equation-codim-1-in-projective-space}
Let $R$ be a Noetherian UFD. Let $Z \subset \mathbf{P}^n_R$ be a
closed subscheme
which has no embedded points such that every irreducible component
of $Z$ has codimension $1$ in $\mathbf{P}^n_R$.
Then the ideal $I(Z) \subset R[T_0, \ldots, T_n]$ corresponding
to $Z$ is principal.
\end{lemma}

\begin{proof}
Observe that the local rings of $X = \mathbf{P}^n_R$ are
UFDs because $X$ is covered by affine pieces isomorphic
to $\mathbf{A}^n_R$ and $R[x_1, \ldots, x_n]$ is a UFD
(Algebra, Lemma \ref{algebra-lemma-polynomial-ring-UFD}).
Thus $Z$ is an effective Cartier divisor by
Lemma \ref{lemma-codimension-1-is-effective-Cartier}.
Let $\mathcal{I} \subset \mathcal{O}_X$ be the quasi-coherent
sheaf of ideals corresponding to $Z$.
Choose an isomorphism $\mathcal{O}(m) \to \mathcal{I}$
for some $m \in \mathbf{Z}$, see
Lemma \ref{lemma-Pic-projective-space-UFD}.
Then the composition
$$
\mathcal{O}_X(m) \to \mathcal{I} \to \mathcal{O}_X
$$
is nonzero. We conclude that $m \leq 0$ and that the corresponding
section of $\mathcal{O}_X(m)^{\otimes -1} = \mathcal{O}_X(-m)$
is given by some $F \in R[T_0, \ldots, T_n]$ of degree $-m$, see
Cohomology of Schemes, Lemma
\ref{coherent-lemma-cohomology-projective-space-over-ring}.
Thus on the $i$th standard open $U_i = D_+(T_i)$ the
closed subscheme $Z \cap U_i$ is cut out by the ideal
$$
(F(T_0/T_i, \ldots, T_n/T_i)) \subset R[T_0/T_i, \ldots, T_n/T_i]
$$
Thus the homogeneous elements of the graded ideal
$I(Z) = \Ker(R[T_0, \ldots, T_n] \to \bigoplus \Gamma(\mathcal{O}_Z(m)))$
is the set of homogeneous polynomials $G$ such that
$$
G(T_0/T_i, \ldots, T_n/T_i) \in (F(T_0/T_i, \ldots, T_n/T_i))
$$
for $i = 0, \ldots, n$. Clearing denominators, we see there exist
$e_i \geq 0$ such that
$$
T_i^{e_i}G \in (F)
$$
for $i = 0, \ldots, n$. As $R$ is a UFD, so is $R[T_0, \ldots, T_n]$.
Then $F | T_0^{e_0}G$ and $F | T_1^{e_1}G$ implies $F | G$ as
$T_0^{e_0}$ and $T_1^{e_1}$ have no factor in common. Thus $I(Z) = (F)$.
\end{proof}

\noindent
In case the closed subscheme is locally cut out by finitely many
equations we can define it by a finite type ideal sheaf of
$\mathcal{A}$.

\begin{lemma}
\label{lemma-closed-subscheme-proj-finite}
Let $S$ be a quasi-compact and quasi-separated scheme.
Let $\mathcal{A}$ be a quasi-coherent graded $\mathcal{O}_S$-algebra. Let
$p : X = \underline{\text{Proj}}_S(\mathcal{A}) \to S$ be the relative
Proj of $\mathcal{A}$. Let $i : Z \to X$ be a closed subscheme.
If $p$ is quasi-compact and $i$ of finite presentation, then there exists
a $d > 0$ and a quasi-coherent finite type $\mathcal{O}_S$-submodule
$\mathcal{F} \subset \mathcal{A}_d$ such that
$Z = \underline{\text{Proj}}_S(\mathcal{A}/\mathcal{F}\mathcal{A})$.
\end{lemma}

\begin{proof}
By Lemma \ref{lemma-closed-subscheme-proj} we know there exists a
quasi-coherent graded sheaf of ideals $\mathcal{I} \subset \mathcal{A}$
such that $Z = \underline{\text{Proj}}(\mathcal{A}/\mathcal{I})$.
Since $S$ is quasi-compact we can choose a finite affine open covering
$S = U_1 \cup \ldots \cup U_n$. Say $U_i = \Spec(R_i)$. Let
$\mathcal{A}|_{U_i}$ correspond to the graded $R_i$-algebra $A_i$ and
$\mathcal{I}|_{U_i}$ to the graded ideal $I_i \subset A_i$. Note that
$p^{-1}(U_i) = \text{Proj}(A_i)$ as schemes over $R_i$.
Since $p$ is quasi-compact we can choose finitely many homogeneous
elements $f_{i, j} \in A_{i, +}$ such that
$p^{-1}(U_i) = \bigcup_j D_{+}(f_{i, j})$.
The condition on $Z \to X$ means that the ideal sheaf of $Z$ in
$\mathcal{O}_X$ is of finite type, see
Morphisms, Lemma \ref{morphisms-lemma-closed-immersion-finite-presentation}.
Hence we can find finitely many homogeneous elements
$h_{i, j, k} \in I_i \cap A_{i, +}$ such that the ideal of
$Z \cap D_{+}(f_{i, j})$ is generated by the elements
$h_{i, j, k}/f_{i, j}^{e_{i, j, k}}$. Choose $d > 0$ to be a common multiple
of all the integers $\deg(f_{i, j})$ and $\deg(h_{i, j, k})$.
By Properties, Lemma \ref{properties-lemma-quasi-coherent-colimit-finite-type}
there exists a finite type quasi-coherent $\mathcal{F} \subset \mathcal{I}_d$
such that all the local sections
$$
h_{i, j, k}f_{i, j}^{(d - \deg(h_{i, j, k}))/\deg(f_{i, j})}
$$
are sections of $\mathcal{F}$. By construction $\mathcal{F}$ is a solution.
\end{proof}

\noindent
The following version of Lemma \ref{lemma-closed-subscheme-proj-finite}
will be used in the proof of
Lemma \ref{lemma-composition-admissible-blowups}.

\begin{lemma}
\label{lemma-closed-subscheme-proj-finite-type}
Let $S$ be a quasi-compact and quasi-separated scheme.
Let $\mathcal{A}$ be a quasi-coherent graded $\mathcal{O}_S$-algebra.
Let $p : X = \underline{\text{Proj}}_S(\mathcal{A}) \to S$ be the relative
Proj of $\mathcal{A}$. Let $i : Z \to X$ be a closed subscheme.
Let $U \subset S$ be an open. Assume that
\begin{enumerate}
\item $p$ is quasi-compact,
\item $i$ of finite presentation,
\item $U \cap p(i(Z)) = \emptyset$,
\item $U$ is quasi-compact,
\item $\mathcal{A}_n$ is a finite type $\mathcal{O}_S$-module for all $n$.
\end{enumerate}
Then there exists a $d > 0$ and a quasi-coherent finite type
$\mathcal{O}_S$-submodule $\mathcal{F} \subset \mathcal{A}_d$ with (a)
$Z = \underline{\text{Proj}}_S(\mathcal{A}/\mathcal{F}\mathcal{A})$
and (b) the support of $\mathcal{A}_d/\mathcal{F}$ is disjoint from $U$.
\end{lemma}

\begin{proof}
Let $\mathcal{I} \subset \mathcal{A}$ be the sheaf of quasi-coherent
graded ideals constructed in Lemma \ref{lemma-closed-subscheme-proj}.
Let $U_i$, $R_i$, $A_i$, $I_i$, $f_{i, j}$, $h_{i, j, k}$, and $d$
be as constructed in the proof of
Lemma \ref{lemma-closed-subscheme-proj-finite}.
Since $U \cap p(i(Z)) = \emptyset$ we see that
$\mathcal{I}_d|_U = \mathcal{A}_d|_U$ (by our construction of
$\mathcal{I}$ as a kernel). Since $U$ is quasi-compact we
can choose a finite affine open covering $U = W_1 \cup \ldots \cup W_m$.
Since $\mathcal{A}_d$ is of finite type we can find finitely many sections
$g_{t, s} \in \mathcal{A}_d(W_t)$ which generate
$\mathcal{A}_d|_{W_t} = \mathcal{I}_d|_{W_t}$
as an $\mathcal{O}_{W_t}$-module. To finish the proof, note that by
Properties, Lemma \ref{properties-lemma-quasi-coherent-colimit-finite-type}
there exists a finite type $\mathcal{F} \subset \mathcal{I}_d$
such that all the local sections
$$
h_{i, j, k}f_{i, j}^{(d - \deg(h_{i, j, k}))/\deg(f_{i, j})}
\quad\text{and}\quad
g_{t, s}
$$
are sections of $\mathcal{F}$. By construction $\mathcal{F}$ is a solution.
\end{proof}

\begin{lemma}
\label{lemma-conormal-sheaf-section-projective-bundle}
Let $X$ be a scheme. Let $\mathcal{E}$ be a quasi-coherent
$\mathcal{O}_X$-module. There is a bijection
$$
\left\{
\begin{matrix}
\text{sections }\sigma\text{ of the } \\
\text{morphism } \mathbf{P}(\mathcal{E}) \to X
\end{matrix}
\right\}
\leftrightarrow
\left\{
\begin{matrix}
\text{surjections }\mathcal{E} \to \mathcal{L}\text{ where} \\
\mathcal{L}\text{ is an invertible }\mathcal{O}_X\text{-module}
\end{matrix}
\right\}
$$
In this case $\sigma$ is a closed immersion and there is a canonical
isomorphism
$$
\Ker(\mathcal{E} \to \mathcal{L})
\otimes_{\mathcal{O}_X} \mathcal{L}^{\otimes -1}
\longrightarrow
\mathcal{C}_{\sigma(X)/\mathbf{P}(\mathcal{E})}
$$
Both the bijection and isomorphism are compatible with base change.
\end{lemma}

\begin{proof}
Recall that $\pi : \mathbf{P}(\mathcal{E}) \to X$ is the relative proj of the
symmetric algebra on $\mathcal{E}$, see
Constructions, Definition \ref{constructions-definition-projective-bundle}.
Hence the descriptions of sections $\sigma$ follows immediately from
the description of the functor of points of $\mathbf{P}(\mathcal{E})$
in Constructions, Lemma \ref{constructions-lemma-apply-relative}.
Since $\pi$ is separated, any section is a closed immersion
(Constructions, Lemma \ref{constructions-lemma-relative-proj-separated} and
Schemes, Lemma \ref{schemes-lemma-section-immersion}).
Let $U \subset X$ be an affine open and $k \in \mathcal{E}(U)$ and
$s \in \mathcal{E}(U)$ be local sections such that $k$ maps to
zero in $\mathcal{L}$ and $s$ maps to a generator $\overline{s}$
of $\mathcal{L}$.
Then $f = k/s$ is a section of $\mathcal{O}_{\mathbf{P}(\mathcal{E})}$
defined in an open neighbourhood $D_+(s)$ of $s(U)$ in $\pi^{-1}(U)$.
Moreover, since $k$ maps to zero in $\mathcal{L}$ we see that
$f$ is a section of the ideal sheaf of $s(U)$ in $\pi^{-1}(U)$.
Thus we can take the image $\overline{f}$ of $f$ in
$\mathcal{C}_{\sigma(X)/\mathbf{P}(\mathcal{E})}(U)$.
We claim (1) that the image $\overline{f}$ depends only on the
sections $k$ and $\overline{s}$ and not on the choice of $s$
and (2) that we get an isomorphism over $U$ in this manner (see below).
However, once (1) and (2) are established, we see that
the construction is compatible with base change by $U' \to U$
where $U'$ is affine, which proves that these local maps glue
and are compatible with arbitrary base change.

\medskip\noindent
To prove (1) and (2) we make explicit what is going on.
Namely, say $U = \Spec(A)$ and say $\mathcal{E} \to \mathcal{L}$
corresponds to the map of $A$-modules $M \to N$. Then
$k \in K = \Ker(M \to N)$ and $s \in M$ maps to a generator $\overline{s}$
of $N$. Hence $M = K \oplus A s$. Thus
$$
\text{Sym}(M) = \text{Sym}(K)[s]
$$
Consider the identification $\text{Sym}(K) \to \text{Sym}(M)_{(s)}$
via the rule $g \mapsto g/s^n$ for $g \in \text{Sym}^n(K)$.
This gives an isomorphism $D_+(s) = \Spec(\text{Sym}(K))$ such
that $\sigma$ corresponds to the ring map $\text{Sym}(K) \to A$
mapping $K$ to zero. Via this isomorphism we see that the quasi-coherent
module corresponding to $K$ is identified with
$\mathcal{C}_{\sigma(U)/D_+(s)}$ proving (2).
Finally, suppose that $s' = k' + s$ for some $k' \in K$.
Then
$$
k/s' = (k/s) (s/s') = (k/s) (s'/s)^{-1} = (k/s) (1 + k'/s)^{-1}
$$
in an open neighbourhood of $\sigma(U)$ in $D_+(s)$. Thus we see that
$s'/s$ restricts to $1$ on $\sigma(U)$ and we see that $k/s'$ maps to
the same element of the conormal sheaf as does $k/s$ thereby proving (1).
\end{proof}





